\documentclass[aspectratio=169]{beamer}
\usepackage[utf8]{inputenc}
\usepackage{lmodern}
\usepackage[T1]{fontenc}
\usepackage{listings}
\usepackage{tikz}
\usetikzlibrary{arrows.meta, positioning, shapes.geometric, calc}

% Set path to find the theme
\makeatletter
\def\input@path{{../BeamerTemplate/theme/}}
\makeatother

\usetheme{CleanEasy}
\input{../BeamerTemplate/configs/configs}

% Listings configuration
\lstset{
    basicstyle=\ttfamily\footnotesize,
    keywordstyle=\color{blue},
    commentstyle=\color{gray},
    stringstyle=\color{red},
    showstringspaces=false,
    breaklines=true,
    frame=single,
    numbers=left,
    numberstyle=\tiny\color{gray}
}

\title{Searching Algorithms}
\subtitle{Find Elements Efficiently Within Collections}
\author{Minseok Jeon}
\institute{DGIST}
\date{\today}

\begin{document}

\begin{frame}
  \titlepage
\end{frame}

\begin{frame}{Table of Contents}
  \tableofcontents
\end{frame}

\section{Introduction}

\begin{frame}{What is Searching?}
  \textbf{Searching}: Finding a specific element in a collection

  \vspace{0.3cm}
  \textbf{Why Searching Matters:}
  \begin{itemize}
    \item Fundamental operation in computer science
    \item Used in databases, search engines, file systems
    \item Performance-critical in many applications
    \item Basis for more complex algorithms
  \end{itemize}

  \vspace{0.3cm}
  \textbf{Key Questions:}
  \begin{itemize}
    \item Is the data sorted?
    \item How large is the dataset?
    \item How frequently do we search?
    \item Can we preprocess the data?
  \end{itemize}
\end{frame}

\begin{frame}{Classification of Search Algorithms}
  \textbf{Based on Data Structure:}
  \begin{itemize}
    \item \textbf{Array-based}: Linear, Binary, Interpolation, Jump
    \item \textbf{Tree-based}: BST, AVL, Red-Black Trees
    \item \textbf{Hash-based}: Hash Tables, Hash Maps
  \end{itemize}

  \vspace{0.3cm}
  \textbf{Based on Preconditions:}
  \begin{itemize}
    \item \textbf{Unsorted data}: Linear Search, Hash Tables
    \item \textbf{Sorted data}: Binary Search, Interpolation, Jump
    \item \textbf{Special structures}: Tree Search
  \end{itemize}

  \vspace{0.3cm}
  \textbf{Based on Complexity:}
  \begin{itemize}
    \item $O(1)$: Hash Table (average)
    \item $O(\log \log n)$: Interpolation (average, uniform data)
    \item $O(\log n)$: Binary Search, BST
    \item $O(\sqrt{n})$: Jump Search
    \item $O(n)$: Linear Search
  \end{itemize}
\end{frame}

\section{Linear vs Binary Search}

\begin{frame}{Linear Search: Overview}
  \textbf{Sequential Search Through Elements}

  \vspace{0.3cm}
  \textbf{Algorithm:}
  \begin{enumerate}
    \item Start from first element
    \item Check each element sequentially
    \item Return index when found
    \item Return -1 if not found
  \end{enumerate}

  \vspace{0.3cm}
  \textbf{Characteristics:}
  \begin{itemize}
    \item \textbf{Time}: $O(n)$
    \item \textbf{Space}: $O(1)$
    \item \textbf{Requirement}: None (works on unsorted data)
  \end{itemize}

  \vspace{0.3cm}
  \textbf{When to Use:}
  \begin{itemize}
    \item Small datasets ($n < 100$)
    \item Unsorted data
    \item Simple implementation needed
    \item One-time search
  \end{itemize}
\end{frame}

\begin{frame}[fragile]{Linear Search: Implementation}
\begin{lstlisting}[language=Python, basicstyle=\ttfamily\small]
def linear_search(arr, target):
    """Search for target in array"""
    for i in range(len(arr)):
        if arr[i] == target:
            return i  # Return index
    return -1  # Not found

# Example
arr = [64, 34, 25, 12, 22, 11, 90]
result = linear_search(arr, 22)  # Returns 4
\end{lstlisting}

\vspace{0.3cm}
\textbf{Analysis:}
\begin{itemize}
  \item \textbf{Best case}: $O(1)$ - element at first position
  \item \textbf{Average case}: $O(n/2) = O(n)$ - element in middle
  \item \textbf{Worst case}: $O(n)$ - element at end or not present
\end{itemize}
\end{frame}

\begin{frame}{Binary Search: Overview}
  \textbf{Divide and Conquer on Sorted Array}

  \vspace{0.3cm}
  \textbf{Algorithm:}
  \begin{enumerate}
    \item Find middle element
    \item Compare with target
    \item If equal: found!
    \item If target $<$ middle: search left half
    \item If target $>$ middle: search right half
    \item Repeat until found or exhausted
  \end{enumerate}

  \vspace{0.3cm}
  \textbf{Characteristics:}
  \begin{itemize}
    \item \textbf{Time}: $O(\log n)$
    \item \textbf{Space}: $O(1)$ iterative, $O(\log n)$ recursive
    \item \textbf{Requirement}: Array must be sorted
  \end{itemize}

  \vspace{0.3cm}
  \textbf{Key Advantage:}
  \begin{itemize}
    \item Eliminates half of search space in each step
    \item For 1 million elements: $\log_2(1,000,000) \approx 20$ comparisons
  \end{itemize}
\end{frame}

\begin{frame}[fragile]{Binary Search: Iterative Implementation}
\begin{lstlisting}[language=Python, basicstyle=\ttfamily\tiny]
def binary_search(arr, target):
    """Search for target in sorted array"""
    left, right = 0, len(arr) - 1

    while left <= right:
        mid = left + (right - left) // 2  # Avoid overflow

        if arr[mid] == target:
            return mid
        elif arr[mid] < target:
            left = mid + 1  # Search right half
        else:
            right = mid - 1  # Search left half

    return -1  # Not found

# Example (array must be sorted!)
arr = [11, 12, 22, 25, 34, 64, 90]
result = binary_search(arr, 22)  # Returns 2
\end{lstlisting}

\textbf{Note:} \texttt{mid = left + (right - left) // 2} avoids integer overflow
\end{frame}

\begin{frame}[fragile]{Binary Search: Recursive Implementation}
\begin{lstlisting}[language=Python, basicstyle=\ttfamily\small]
def binary_search_recursive(arr, target, left, right):
    if left > right:
        return -1

    mid = left + (right - left) // 2

    if arr[mid] == target:
        return mid
    elif arr[mid] < target:
        return binary_search_recursive(arr, target,
                                       mid + 1, right)
    else:
        return binary_search_recursive(arr, target,
                                       left, mid - 1)

# Usage
arr = [11, 12, 22, 25, 34, 64, 90]
result = binary_search_recursive(arr, 22, 0, len(arr) - 1)
\end{lstlisting}
\end{frame}

\begin{frame}{Binary Search: Visual Example}
  \textbf{Search for 22 in [11, 12, 22, 25, 34, 64, 90]}

  \vspace{0.3cm}
  \textbf{Step 1:} left=0, right=6, mid=3
  \begin{center}
  [11, 12, 22, {\color{red}25}, 34, 64, 90]

  arr[3]=25 $>$ 22, search left half
  \end{center}

  \vspace{0.2cm}
  \textbf{Step 2:} left=0, right=2, mid=1
  \begin{center}
  [11, {\color{red}12}, 22]

  arr[1]=12 $<$ 22, search right half
  \end{center}

  \vspace{0.2cm}
  \textbf{Step 3:} left=2, right=2, mid=2
  \begin{center}
  [{\color{green}22}]

  arr[2]=22, found at index 2!
  \end{center}

  \vspace{0.3cm}
  \textbf{Result:} Only 3 comparisons for 7 elements
\end{frame}

\begin{frame}{Linear vs Binary Search: Comparison}
  \begin{center}
  \begin{tabular}{|l|c|c|}
    \hline
    \textbf{Feature} & \textbf{Linear} & \textbf{Binary} \\
    \hline
    Time Complexity & $O(n)$ & $O(\log n)$ \\
    \hline
    Space Complexity & $O(1)$ & $O(1)$ iterative \\
    \hline
    Sorted Required & No & Yes \\
    \hline
    Best For & Small/unsorted & Large/sorted \\
    \hline
    Avg comparisons (n=1000) & 500 & 10 \\
    \hline
    Implementation & Simple & Moderate \\
    \hline
    Linked List Support & Yes & No \\
    \hline
  \end{tabular}
  \end{center}

  \vspace{0.3cm}
  \textbf{Performance Comparison:}
  \begin{itemize}
    \item $n=10$: Linear=5, Binary=3
    \item $n=100$: Linear=50, Binary=7
    \item $n=1,000,000$: Linear=500,000, Binary=20
  \end{itemize}
\end{frame}

\section{Binary Search on Answers (Parametric)}

\begin{frame}{Concept: Search the Answer Space}
  \textbf{Binary Search on Answers (Parametric Search)}

  \vspace{0.3cm}
  \textbf{Idea:}
  \begin{itemize}
    \item Binary search on the \textbf{solution space}, not the array
    \item Find minimum/maximum value satisfying a condition
    \item Requires monotonic property
  \end{itemize}

  \vspace{0.3cm}
  \textbf{Monotonic Property:}
  \begin{itemize}
    \item If value $x$ works, then all values $> x$ also work (or vice versa)
    \item Creates a boundary: [false, false, ..., true, true, ...]
    \item Binary search finds the boundary
  \end{itemize}

  \vspace{0.3cm}
  \textbf{Common Patterns:}
  \begin{itemize}
    \item \textbf{Minimize maximum}: Find smallest value where all constraints satisfied
    \item \textbf{Maximize minimum}: Find largest value where all constraints satisfied
    \item \textbf{Find threshold}: Find boundary where condition changes
  \end{itemize}
\end{frame}

\begin{frame}[fragile]{Template: Binary Search on Answers}
\begin{lstlisting}[language=Python, basicstyle=\ttfamily\small]
def binary_search_answer(condition_func, low, high):
    """
    Find minimum value in [low, high] satisfying condition
    """
    result = -1

    while low <= high:
        mid = low + (high - low) // 2

        if condition_func(mid):
            result = mid
            high = mid - 1  # Try to find smaller
        else:
            low = mid + 1

    return result
\end{lstlisting}

\textbf{Key Points:}
\begin{itemize}
  \item \texttt{condition\_func(x)}: Returns True if x satisfies constraint
  \item Search space: continuous range of values
  \item Answer: boundary where condition becomes true
\end{itemize}
\end{frame}

\begin{frame}[fragile]{Example 1: Integer Square Root}
\begin{lstlisting}[language=Python, basicstyle=\ttfamily\tiny]
def sqrt(x):
    """Find integer square root of x"""
    if x < 2:
        return x

    left, right = 1, x // 2

    while left <= right:
        mid = left + (right - left) // 2
        square = mid * mid

        if square == x:
            return mid
        elif square < x:
            left = mid + 1
        else:
            right = mid - 1

    return right  # Largest integer whose square <= x

# Example: sqrt(8) = 2 (since 2^2 = 4 <= 8 < 3^2 = 9)
print(sqrt(8))   # Output: 2
print(sqrt(16))  # Output: 4
print(sqrt(17))  # Output: 4
\end{lstlisting}

\textbf{Analysis:} Search space is [1, x/2], binary search in $O(\log x)$
\end{frame}

\begin{frame}[fragile]{Example 2: Minimum Ship Capacity}
\textbf{Problem:} Ship packages in D days, find minimum capacity

\begin{lstlisting}[language=Python, basicstyle=\ttfamily\tiny]
def ship_within_days(weights, days):
    """Find minimum ship capacity"""
    def can_ship(capacity):
        """Check if we can ship with given capacity"""
        days_needed = 1
        current_load = 0

        for weight in weights:
            if current_load + weight > capacity:
                days_needed += 1
                current_load = weight
            else:
                current_load += weight

        return days_needed <= days

    # Binary search on capacity
    left = max(weights)  # Min: heaviest package
    right = sum(weights)  # Max: all at once

    while left < right:
        mid = left + (right - left) // 2

        if can_ship(mid):
            right = mid  # Try smaller capacity
        else:
            left = mid + 1

    return left
\end{lstlisting}
\end{frame}

\begin{frame}{Ship Capacity Example}
  \textbf{Input:} weights = [1,2,3,4,5,6,7,8,9,10], days = 5

  \vspace{0.3cm}
  \textbf{Search Space:}
  \begin{itemize}
    \item Minimum capacity: max(weights) = 10
    \item Maximum capacity: sum(weights) = 55
    \item Search range: [10, 55]
  \end{itemize}

  \vspace{0.3cm}
  \textbf{Binary Search Process:}
  \begin{itemize}
    \item Try capacity=32: Can ship in 3 days $\rightarrow$ too large
    \item Try capacity=21: Can ship in 4 days $\rightarrow$ too large
    \item Try capacity=15: Can ship in 5 days $\rightarrow$ works!
    \item Try capacity=12: Cannot ship in 5 days $\rightarrow$ too small
  \end{itemize}

  \vspace{0.3cm}
  \textbf{Answer:} 15

  \textbf{Distribution:} [1,2,3,4,5], [6,7], [8], [9], [10]
\end{frame}

\begin{frame}[fragile]{Example 3: Koko Eating Bananas}
\textbf{Problem:} Finish all banana piles in H hours, minimize eating speed

\begin{lstlisting}[language=Python, basicstyle=\ttfamily\tiny]
def min_eating_speed(piles, h):
    """Find minimum eating speed"""
    def can_finish(speed):
        """Check if can finish with given speed"""
        hours = 0
        for pile in piles:
            hours += (pile + speed - 1) // speed  # Ceiling
        return hours <= h

    left, right = 1, max(piles)

    while left < right:
        mid = left + (right - left) // 2

        if can_finish(mid):
            right = mid
        else:
            left = mid + 1

    return left

# Example: piles=[3,6,7,11], h=8
# Answer: 4 (eat 4 bananas/hour)
\end{lstlisting}
\end{frame}

\section{Search Trees and Hashing}

\begin{frame}{Binary Search Tree (BST)}
  \textbf{Tree-Based Search Structure}

  \vspace{0.3cm}
  \textbf{Properties:}
  \begin{itemize}
    \item Each node has left (smaller) and right (larger) children
    \item Left subtree $<$ node $<$ right subtree
    \item Enables $O(\log n)$ search on average
  \end{itemize}

  \vspace{0.3cm}
  \textbf{Characteristics:}
  \begin{itemize}
    \item \textbf{Search}: $O(\log n)$ average, $O(n)$ worst
    \item \textbf{Insert/Delete}: $O(\log n)$ average
    \item \textbf{Space}: $O(n)$
  \end{itemize}

  \vspace{0.3cm}
  \textbf{Advantages:}
  \begin{itemize}
    \item Dynamic: supports insertions/deletions
    \item Maintains sorted order
    \item Range queries efficient
  \end{itemize}

  \vspace{0.3cm}
  \textbf{Disadvantage:}
  \begin{itemize}
    \item Can degenerate to $O(n)$ if unbalanced
    \item Solution: Self-balancing trees (AVL, Red-Black)
  \end{itemize}
\end{frame}

\begin{frame}[fragile]{BST Search Implementation}
\begin{lstlisting}[language=Python, basicstyle=\ttfamily\small]
class TreeNode:
    def __init__(self, val):
        self.val = val
        self.left = None
        self.right = None

def search_bst(root, target):
    """Search in BST - Recursive"""
    if not root or root.val == target:
        return root

    if target < root.val:
        return search_bst(root.left, target)
    else:
        return search_bst(root.right, target)

def search_bst_iterative(root, target):
    """Search in BST - Iterative"""
    while root and root.val != target:
        if target < root.val:
            root = root.left
        else:
            root = root.right
    return root
\end{lstlisting}
\end{frame}

\begin{frame}{Hash Table Search}
  \textbf{Constant-Time Lookup via Hashing}

  \vspace{0.3cm}
  \textbf{Concept:}
  \begin{itemize}
    \item Map keys to array indices using hash function
    \item Direct access to value via computed index
    \item Handle collisions (chaining or open addressing)
  \end{itemize}

  \vspace{0.3cm}
  \textbf{Characteristics:}
  \begin{itemize}
    \item \textbf{Search}: $O(1)$ average, $O(n)$ worst
    \item \textbf{Insert/Delete}: $O(1)$ average
    \item \textbf{Space}: $O(n)$
  \end{itemize}

  \vspace{0.3cm}
  \textbf{Advantages:}
  \begin{itemize}
    \item Fastest average-case search: $O(1)$
    \item Simple interface (key-value pairs)
    \item Good for large datasets with unique keys
  \end{itemize}

  \vspace{0.3cm}
  \textbf{Disadvantages:}
  \begin{itemize}
    \item No sorted order
    \item Poor worst-case performance
    \item Extra memory for hash table
  \end{itemize}
\end{frame}

\begin{frame}[fragile]{Hash Table Implementation}
\begin{lstlisting}[language=Python, basicstyle=\ttfamily\tiny]
class HashTable:
    def __init__(self, size=10):
        self.size = size
        self.table = [[] for _ in range(size)]

    def _hash(self, key):
        return hash(key) % self.size

    def insert(self, key, value):
        index = self._hash(key)
        for i, (k, v) in enumerate(self.table[index]):
            if k == key:
                self.table[index][i] = (key, value)
                return
        self.table[index].append((key, value))

    def search(self, key):
        index = self._hash(key)
        for k, v in self.table[index]:
            if k == key:
                return v
        return None

# Python's built-in dict is a hash table
hash_table = {"apple": 5, "banana": 3}
print(hash_table["apple"])  # O(1) lookup
\end{lstlisting}
\end{frame}

\begin{frame}{Comparison: Array, BST, Hash Table}
  \begin{center}
  \begin{tabular}{|l|c|c|c|}
    \hline
    \textbf{Structure} & \textbf{Average} & \textbf{Worst} & \textbf{Sorted} \\
    \hline
    Sorted Array + Binary & $O(\log n)$ & $O(\log n)$ & Yes \\
    \hline
    BST & $O(\log n)$ & $O(n)$ & Yes \\
    \hline
    Balanced BST & $O(\log n)$ & $O(\log n)$ & Yes \\
    \hline
    Hash Table & $O(1)$ & $O(n)$ & No \\
    \hline
  \end{tabular}
  \end{center}

  \vspace{0.3cm}
  \textbf{When to Use:}
  \begin{itemize}
    \item \textbf{Sorted Array}: Static data, range queries
    \item \textbf{BST}: Dynamic data, need sorted order
    \item \textbf{Balanced BST}: Guaranteed performance, sorted
    \item \textbf{Hash Table}: Fast lookup, no order needed
  \end{itemize}
\end{frame}

\section{Interpolation and Exponential Search}

\begin{frame}{Interpolation Search}
  \textbf{Position Estimation Based on Value}

  \vspace{0.3cm}
  \textbf{Idea:}
  \begin{itemize}
    \item Like looking up a name in phone book
    \item Estimate position based on value distribution
    \item Works best with uniformly distributed data
  \end{itemize}

  \vspace{0.3cm}
  \textbf{Position Formula:}
  \begin{center}
  $pos = left + \frac{(target - arr[left])}{(arr[right] - arr[left])} \times (right - left)$
  \end{center}

  \vspace{0.3cm}
  \textbf{Characteristics:}
  \begin{itemize}
    \item \textbf{Time}: $O(\log \log n)$ average, $O(n)$ worst
    \item \textbf{Requirement}: Sorted + uniformly distributed
    \item \textbf{Better than binary}: For uniform data
  \end{itemize}

  \vspace{0.3cm}
  \textbf{Example:} Array [10, 20, 30, 40, 50, 60, 70, 80, 90], target=70

  Estimate: $pos \approx 0 + \frac{70-10}{90-10} \times 8 = 6$ (directly finds it!)
\end{frame}

\begin{frame}[fragile]{Interpolation Search Implementation}
\begin{lstlisting}[language=Python, basicstyle=\ttfamily\tiny]
def interpolation_search(arr, target):
    """Search in uniformly distributed sorted array"""
    left, right = 0, len(arr) - 1

    while left <= right and target >= arr[left] and target <= arr[right]:
        if left == right:
            if arr[left] == target:
                return left
            return -1

        # Estimate position
        pos = left + int(((target - arr[left]) /
                         (arr[right] - arr[left])) * (right - left))

        if arr[pos] == target:
            return pos
        elif arr[pos] < target:
            left = pos + 1
        else:
            right = pos - 1

    return -1

# Example: uniformly distributed data
arr = [10, 20, 30, 40, 50, 60, 70, 80, 90]
result = interpolation_search(arr, 70)  # Returns 6
\end{lstlisting}
\end{frame}

\begin{frame}{Exponential Search}
  \textbf{Find Range, Then Binary Search}

  \vspace{0.3cm}
  \textbf{Algorithm:}
  \begin{enumerate}
    \item Check if element at position 0
    \item Find range $[2^{k-1}, 2^k]$ where element exists
    \item Perform binary search in that range
  \end{enumerate}

  \vspace{0.3cm}
  \textbf{Characteristics:}
  \begin{itemize}
    \item \textbf{Time}: $O(\log n)$
    \item \textbf{Best for}: Element near beginning, unbounded arrays
  \end{itemize}

  \vspace{0.3cm}
  \textbf{Why Useful:}
  \begin{itemize}
    \item Don't need to know array size
    \item Better than binary search when target near start
    \item Good for infinite/unbounded lists
  \end{itemize}

  \vspace{0.3cm}
  \textbf{Example:} Search for 10 in sorted array

  Check positions: 1, 2, 4, 8, 16 $\rightarrow$ found range [8, 16]

  Binary search in [8, 16] $\rightarrow$ find exact position
\end{frame}

\begin{frame}[fragile]{Exponential Search Implementation}
\begin{lstlisting}[language=Python, basicstyle=\ttfamily\tiny]
def exponential_search(arr, target):
    """Search in sorted array, efficient when target near start"""
    n = len(arr)

    # If target at first position
    if arr[0] == target:
        return 0

    # Find range for binary search
    i = 1
    while i < n and arr[i] <= target:
        i *= 2

    # Binary search in range [i//2, min(i, n-1)]
    return binary_search_range(arr, target, i // 2, min(i, n - 1))

def binary_search_range(arr, target, left, right):
    while left <= right:
        mid = left + (right - left) // 2

        if arr[mid] == target:
            return mid
        elif arr[mid] < target:
            left = mid + 1
        else:
            right = mid - 1

    return -1
\end{lstlisting}
\end{frame}

\begin{frame}{Jump Search}
  \textbf{Jump Fixed Steps, Then Linear}

  \vspace{0.3cm}
  \textbf{Algorithm:}
  \begin{enumerate}
    \item Jump ahead by $\sqrt{n}$ steps
    \item Find block containing target
    \item Linear search within block
  \end{enumerate}

  \vspace{0.3cm}
  \textbf{Characteristics:}
  \begin{itemize}
    \item \textbf{Time}: $O(\sqrt{n})$
    \item \textbf{Jump size}: Optimal is $\sqrt{n}$
    \item \textbf{Space}: $O(1)$
  \end{itemize}

  \vspace{0.3cm}
  \textbf{Why Use:}
  \begin{itemize}
    \item Simpler than binary search
    \item Better than linear search
    \item Good for systems where backward jumps expensive
  \end{itemize}

  \vspace{0.3cm}
  \textbf{Trade-off:}
  \begin{itemize}
    \item Slower than binary search: $O(\sqrt{n})$ vs $O(\log n)$
    \item Simpler implementation
    \item Better cache performance (forward-only)
  \end{itemize}
\end{frame}

\begin{frame}[fragile]{Jump Search Implementation}
\begin{lstlisting}[language=Python, basicstyle=\ttfamily\small]
import math

def jump_search(arr, target):
    """Jump search in sorted array"""
    n = len(arr)
    step = int(math.sqrt(n))
    prev = 0

    # Jump to find block
    while arr[min(step, n) - 1] < target:
        prev = step
        step += int(math.sqrt(n))
        if prev >= n:
            return -1

    # Linear search in block
    while arr[prev] < target:
        prev += 1
        if prev == min(step, n):
            return -1

    if arr[prev] == target:
        return prev

    return -1
\end{lstlisting}
\end{frame}

\section{Complexity and Preconditions}

\begin{frame}{Time Complexity Summary}
  \begin{center}
  \begin{tabular}{|l|c|c|c|c|}
    \hline
    \textbf{Algorithm} & \textbf{Best} & \textbf{Avg} & \textbf{Worst} & \textbf{Space} \\
    \hline
    Linear & $O(1)$ & $O(n)$ & $O(n)$ & $O(1)$ \\
    \hline
    Binary & $O(1)$ & $O(\log n)$ & $O(\log n)$ & $O(1)$ \\
    \hline
    Interpolation & $O(1)$ & $O(\log \log n)$ & $O(n)$ & $O(1)$ \\
    \hline
    Exponential & $O(1)$ & $O(\log n)$ & $O(\log n)$ & $O(1)$ \\
    \hline
    Jump & $O(1)$ & $O(\sqrt{n})$ & $O(\sqrt{n})$ & $O(1)$ \\
    \hline
    BST & $O(\log n)$ & $O(\log n)$ & $O(n)$ & $O(n)$ \\
    \hline
    Hash Table & $O(1)$ & $O(1)$ & $O(n)$ & $O(n)$ \\
    \hline
  \end{tabular}
  \end{center}

  \vspace{0.3cm}
  \textbf{Performance Comparison (n=1,000,000):}
  \begin{itemize}
    \item Linear: up to 1,000,000 comparisons
    \item Jump: $\sqrt{1,000,000} = 1,000$ comparisons
    \item Binary: $\log_2(1,000,000) \approx 20$ comparisons
    \item Interpolation: $\log \log(1,000,000) \approx 4$ comparisons
    \item Hash: 1 comparison (average)
  \end{itemize}
\end{frame}

\begin{frame}[fragile]{Preconditions: Binary Search}
\textbf{MUST be sorted!}

\begin{lstlisting}[language=Python, basicstyle=\ttfamily\small]
# Correct usage
arr = [1, 3, 5, 7, 9, 11, 13]  # Sorted
result = binary_search(arr, 7)  # Works correctly

# Incorrect usage
arr = [3, 1, 5, 9, 7, 11, 13]  # NOT sorted
result = binary_search(arr, 7)  # May fail!

# Check if sorted
def is_sorted(arr):
    return all(arr[i] <= arr[i+1]
               for i in range(len(arr)-1))

# Safe usage
if is_sorted(arr):
    result = binary_search(arr, target)
else:
    arr.sort()  # Sort first
    result = binary_search(arr, target)
\end{lstlisting}
\end{frame}

\begin{frame}[fragile]{Preconditions: Other Algorithms}
\textbf{Interpolation Search:}
\begin{lstlisting}[language=Python, basicstyle=\ttfamily\small]
# MUST be sorted AND uniformly distributed
arr = [10, 20, 30, 40, 50]        # Good: uniform
arr = [1, 2, 3, 100, 1000]        # Bad: not uniform
\end{lstlisting}

\vspace{0.3cm}
\textbf{Hash Table:}
\begin{lstlisting}[language=Python, basicstyle=\ttfamily\small]
# Keys must be hashable (immutable)
hash_table[42] = "value"           # OK: int
hash_table["key"] = "value"        # OK: str
hash_table[(1, 2)] = "value"       # OK: tuple
hash_table[[1, 2]] = "value"       # ERROR: list!
\end{lstlisting}

\vspace{0.3cm}
\textbf{BST:}
\begin{lstlisting}[language=Python, basicstyle=\ttfamily\small]
# Elements must be comparable
bst.insert(5)    # OK: int
bst.insert("a")  # ERROR if tree has ints!
\end{lstlisting}
\end{frame}

\begin{frame}{Choosing the Right Algorithm}
  \textbf{Decision Factors:}

  \vspace{0.3cm}
  \textbf{1. Data Characteristics:}
  \begin{itemize}
    \item Sorted? $\rightarrow$ Binary, Interpolation, Jump
    \item Uniformly distributed? $\rightarrow$ Interpolation
    \item Small size ($< 100$)? $\rightarrow$ Linear
    \item Large size? $\rightarrow$ Binary or Hash
  \end{itemize}

  \vspace{0.3cm}
  \textbf{2. Operation Frequency:}
  \begin{itemize}
    \item Many searches? $\rightarrow$ Hash Table or BST
    \item One-time search? $\rightarrow$ Linear
    \item Frequent insertions/deletions? $\rightarrow$ Hash or BST
  \end{itemize}

  \vspace{0.3cm}
  \textbf{3. Requirements:}
  \begin{itemize}
    \item Need sorted order? $\rightarrow$ BST or Sorted Array
    \item Range queries? $\rightarrow$ BST
    \item Fastest possible? $\rightarrow$ Hash Table
    \item Memory constrained? $\rightarrow$ In-place algorithms
  \end{itemize}
\end{frame}

\section{Applications and Patterns}

\begin{frame}[fragile]{Pattern 1: Finding Boundaries}
\textbf{Find first and last occurrence in sorted array}

\begin{lstlisting}[language=Python, basicstyle=\ttfamily\tiny]
def find_first_occurrence(arr, target):
    """Find first occurrence of target"""
    left, right = 0, len(arr) - 1
    result = -1

    while left <= right:
        mid = left + (right - left) // 2

        if arr[mid] == target:
            result = mid
            right = mid - 1  # Continue searching left
        elif arr[mid] < target:
            left = mid + 1
        else:
            right = mid - 1

    return result

def find_last_occurrence(arr, target):
    """Find last occurrence of target"""
    left, right = 0, len(arr) - 1
    result = -1

    while left <= right:
        mid = left + (right - left) // 2

        if arr[mid] == target:
            result = mid
            left = mid + 1  # Continue searching right
        elif arr[mid] < target:
            left = mid + 1
        else:
            right = mid - 1

    return result
\end{lstlisting}
\end{frame}

\begin{frame}[fragile]{Pattern 2: Rotated Array Search}
\textbf{Search in rotated sorted array}

\begin{lstlisting}[language=Python, basicstyle=\ttfamily\tiny]
def search_rotated(arr, target):
    """
    Search in rotated sorted array
    Example: [4,5,6,7,0,1,2] (rotated at index 4)
    """
    left, right = 0, len(arr) - 1

    while left <= right:
        mid = left + (right - left) // 2

        if arr[mid] == target:
            return mid

        # Determine which half is sorted
        if arr[left] <= arr[mid]:
            # Left half is sorted
            if arr[left] <= target < arr[mid]:
                right = mid - 1
            else:
                left = mid + 1
        else:
            # Right half is sorted
            if arr[mid] < target <= arr[right]:
                left = mid + 1
            else:
                right = mid - 1

    return -1
\end{lstlisting}
\end{frame}

\begin{frame}[fragile]{Pattern 3: Peak Finding}
\begin{lstlisting}[language=Python, basicstyle=\ttfamily\small]
def find_peak_element(arr):
    """
    Find a peak element (greater than neighbors)
    """
    left, right = 0, len(arr) - 1

    while left < right:
        mid = left + (right - left) // 2

        if arr[mid] < arr[mid + 1]:
            left = mid + 1  # Peak is on right
        else:
            right = mid  # Peak is on left or at mid

    return left

# Example: [1, 3, 20, 4, 1, 0]
# Peak at index 2 (value 20)
\end{lstlisting}

\textbf{Key Insight:} Always move towards higher values
\end{frame}

\begin{frame}[fragile]{Pattern 4: 2D Matrix Search}
\begin{lstlisting}[language=Python, basicstyle=\ttfamily\tiny]
def search_matrix(matrix, target):
    """
    Search in row-wise and column-wise sorted matrix
    Example: [[1,4,7,11],
              [2,5,8,12],
              [3,6,9,16],
              [10,13,14,17]]
    """
    if not matrix or not matrix[0]:
        return False

    rows, cols = len(matrix), len(matrix[0])

    # Start from top-right corner
    row, col = 0, cols - 1

    while row < rows and col >= 0:
        if matrix[row][col] == target:
            return True
        elif matrix[row][col] > target:
            col -= 1  # Move left
        else:
            row += 1  # Move down

    return False
\end{lstlisting}

\textbf{Time}: $O(m + n)$, \textbf{Strategy}: Eliminate row or column each step
\end{frame}

\begin{frame}[fragile]{Pattern 5: Finding Missing Number}
\begin{lstlisting}[language=Python, basicstyle=\ttfamily\small]
def find_missing(arr):
    """
    Find missing number in [0, n]
    Example: [0,1,3,4,5] -> missing 2
    """
    left, right = 0, len(arr) - 1

    while left <= right:
        mid = left + (right - left) // 2

        if arr[mid] == mid:
            left = mid + 1  # Missing is on right
        else:
            right = mid - 1  # Missing is on left

    return left

# Example usage
arr = [0, 1, 3, 4, 5, 6]
print(find_missing(arr))  # Output: 2
\end{lstlisting}

\textbf{Key}: If arr[i] == i, no missing number up to i
\end{frame}

\begin{frame}{Common Search Patterns Summary}
  \textbf{1. Find Exact Match:}
  \begin{itemize}
    \item Standard binary search
    \item Return index or -1
  \end{itemize}

  \vspace{0.2cm}
  \textbf{2. Find First/Last Occurrence:}
  \begin{itemize}
    \item Binary search with boundary tracking
    \item Continue searching after finding match
  \end{itemize}

  \vspace{0.2cm}
  \textbf{3. Find Insertion Position:}
  \begin{itemize}
    \item Binary search returning left pointer
    \item Used for maintaining sorted order
  \end{itemize}

  \vspace{0.2cm}
  \textbf{4. Search in Rotated Array:}
  \begin{itemize}
    \item Identify which half is sorted
    \item Search appropriate half
  \end{itemize}

  \vspace{0.2cm}
  \textbf{5. Find Peak/Valley:}
  \begin{itemize}
    \item Compare with neighbors
    \item Move towards increasing values
  \end{itemize}

  \vspace{0.2cm}
  \textbf{6. 2D Search:}
  \begin{itemize}
    \item Start from corner
    \item Eliminate row or column each step
  \end{itemize}

  \vspace{0.2cm}
  \textbf{7. Binary Search on Answer:}
  \begin{itemize}
    \item Search solution space
    \item Check monotonic condition
  \end{itemize}
\end{frame}

\section{Summary}

\begin{frame}{Key Takeaways}
  \textbf{Search Algorithm Categories:}
  \begin{itemize}
    \item \textbf{Linear Search}: $O(n)$, works on any data
    \item \textbf{Binary Search}: $O(\log n)$, requires sorted data
    \item \textbf{Interpolation}: $O(\log \log n)$, uniform distribution
    \item \textbf{Hash Table}: $O(1)$ average, no order
    \item \textbf{BST}: $O(\log n)$, dynamic with sorted order
  \end{itemize}

  \vspace{0.3cm}
  \textbf{Key Concepts:}
  \begin{itemize}
    \item Binary search eliminates half each step
    \item Parametric search: search answer space
    \item Hash tables trade space for speed
    \item Preconditions matter (sorted, uniform, etc.)
  \end{itemize}

  \vspace{0.3cm}
  \textbf{Choosing Algorithm:}
  \begin{itemize}
    \item Consider: data size, sorted, distribution, frequency
    \item Small/unsorted: Linear
    \item Large/sorted: Binary
    \item Frequent searches: Hash or BST
  \end{itemize}
\end{frame}

\begin{frame}{Practical Recommendations}
  \textbf{Use Built-in Functions:}
  \begin{itemize}
    \item Python: \texttt{list.index()}, \texttt{bisect} module, \texttt{dict}
    \item Java: \texttt{Arrays.binarySearch()}, \texttt{HashMap}
    \item C++: \texttt{std::binary\_search()}, \texttt{std::map}
  \end{itemize}

  \vspace{0.3cm}
  \textbf{Common Mistakes to Avoid:}
  \begin{itemize}
    \item Forgetting to sort before binary search
    \item Integer overflow in mid calculation
    \item Off-by-one errors in loop conditions
    \item Using linear search on large sorted data
  \end{itemize}

  \vspace{0.3cm}
  \textbf{Optimization Tips:}
  \begin{itemize}
    \item Preprocess data (sort, build index)
    \item Cache frequent queries
    \item Use appropriate data structure
    \item Consider trade-offs: time vs space vs complexity
  \end{itemize}
\end{frame}

\begin{frame}{Practice Problems}
  \textbf{Problem 1: Binary Search Variants}
  \begin{itemize}
    \item Implement finding first and last occurrence
    \item Find insertion position for sorted array
  \end{itemize}

  \vspace{0.2cm}
  \textbf{Problem 2: Parametric Search}
  \begin{itemize}
    \item Koko eating bananas (LeetCode 875)
    \item Capacity to ship packages (LeetCode 1011)
    \item Split array largest sum (LeetCode 410)
  \end{itemize}

  \vspace{0.2cm}
  \textbf{Problem 3: Rotated Array}
  \begin{itemize}
    \item Search in rotated sorted array
    \item Find minimum in rotated array
  \end{itemize}

  \vspace{0.2cm}
  \textbf{Problem 4: 2D Search}
  \begin{itemize}
    \item Search 2D matrix (LeetCode 74, 240)
    \item Find peak in 2D array
  \end{itemize}
\end{frame}

\begin{frame}{Resources}
  \textbf{Books:}
  \begin{itemize}
    \item "Introduction to Algorithms" (CLRS) - Chapters 2, 11, 12
    \item "The Algorithm Design Manual" (Skiena)
  \end{itemize}

  \vspace{0.3cm}
  \textbf{Online Resources:}
  \begin{itemize}
    \item Visualizations: visualgo.net
    \item LeetCode Binary Search problems
    \item GeeksforGeeks search algorithm tutorials
  \end{itemize}

  \vspace{0.3cm}
  \textbf{Practice Platforms:}
  \begin{itemize}
    \item LeetCode: Binary Search tag
    \item HackerRank: Search challenges
    \item Codeforces: Binary search problems
  \end{itemize}

  \vspace{0.3cm}
  \textbf{Advanced Topics:}
  \begin{itemize}
    \item Ternary search
    \item Fractional cascading
    \item Suffix arrays and trees
  \end{itemize}
\end{frame}

\end{document}
